% TOP_PROBLEM: {"id":30007591,"problem_number":"LOCAL-30007591","title":"Welfare effects of nudges","category_id":21,"category":{"id":21,"name":"economics","display_name":"Economics","description":"Open economic theory statements, published research agendas, and proposed scoped specifications; question provenance and historical priority are qualified per record.","slug":"economics","order_index":21,"created_at":"2026-10-08T00:00:00Z"},"status":"provenance_pending","external_url":null,"legacy_ids":[],"published":false,"statement":"Estimate the specified welfare effect of a nudge using an explicit welfare criterion and identify which conclusions remain robust to uncertain preferences.","economics":{"original_id":80,"field":"Behavioral and experimental economics","kind":"identification","formalization_basis":"proposed_specification","source_status":"not_verified","priority_status":"not_established","priority_note":"The cited paper already derives welfare conditions in its framework. No exact remaining general open-problem locator or historical priority was verified.","earliest_verified_reference":null,"current_reference":{"authors":["Hunt Allcott","Daniel Cohen","William Morrison","Dmitry Taubinsky"],"title":"When Do \"Nudges\" Increase Welfare?","year":2025,"url":"https://www.aeaweb.org/articles?id=10.1257/aer.20231304","doi":"10.1257/aer.20231304","locator":"Abstract; AER 115(5), pp. 1555–1596. Full paper was not accessible in this audit.","evidence_summary":"The abstract explicitly presents welfare conditions involving the variance of distortions, pass-through, and optimal taxes as results. It does not establish the catalogue wording as an unresolved theorem or verify a remaining general spillover-inclusive agenda."},"formalization_note":"Proposed application-level identification question. The original broad wording mixes theoretical results already supplied by Allcott et al. with unverified extensions; retain with an honest unsupported-status flag."},"statusQualification":"Provenance pending: an explicit published statement of this exact open question was not verified. This is a catalogue-authored candidate specification, not a certified open conjecture. Proposed application-level identification question. The original broad wording mixes theoretical results already supplied by Allcott et al. with unverified extensions; retain with an honest unsupported-status flag.","statusReviewedAt":"2026-10-08"}
% FIRST_POSTED: 2026-10-08
% ENTRY_KIND: statement_only
% !TeX program = lualatex
\documentclass[11pt]{article}
\usepackage{fontspec}
\usepackage[a4paper,margin=25mm]{geometry}
\usepackage{amsmath,amssymb,mathtools}
\usepackage{xurl}
\usepackage[hidelinks]{hyperref}
\newcommand{\sref}[2]{\hyperlink{source:#1:#2}{[S#2]}}
\setlength{\emergencystretch}{3em}
\begin{document}
% problemId:30007591
\section*{Welfare effects of nudges}\label{30007591}
\subsection{Definitions and mathematical statement}
Consider consumers \(i\) choosing quantities \(q_i\) of a good. Let \(v_i(q)\) be experienced monetary-equivalent benefit, \(b_i(q,n)\) a choice distortion under nudge \(n\in\{0,1\}\), and \(p(n,t)\) the equilibrium price when a tax \(t\) is fixed. A specified behavioral demand model is
\[q_i(n,t)\in\arg\max_{q\geq0}\{v_i(q)+b_i(q,n)-(p(n,t)+t)q\}.\]
Define welfare, including production costs \(C\), aggregate external damage \(H\), and real implementation cost \(K\), by
\[W(n,t)=\mathbb E[v_i(q_i(n,t))]-C(Q(n,t))-H(Q(n,t))-K(n),\qquad Q(n,t)=\mathbb E[q_i(n,t)].\]
Transfers from prices and tax receipts cancel only under the chosen welfare weights and incidence assumptions. The target estimand is \(\Delta W(t)=W(1,t)-W(0,t)\), or the corresponding difference with taxes reoptimized under each nudge. Determine which experiments and market data identify experienced benefits, heterogeneous distortions, price responses, and spillovers sufficiently to sign \(\Delta W\) in a specified policy application. A mean change in purchases does not by itself identify welfare, and observed choice utility need not equal experienced benefit. Report sensitivity to benefit measurement, welfare weights, unobserved heterogeneity, and equilibrium assumptions. This is an application-level identification proposal; it does not assert that the general welfare conditions studied in the cited paper remain mathematically unsolved.

\par\textbf{Formulation and scope.} Proposed application-level identification question. The original broad wording mixes theoretical results already supplied by Allcott et al. with unverified extensions; retain with an honest unsupported-status flag.

\par\textbf{Open-question provenance.} An explicit published open-question statement was not verified. Sources below are supporting literature and must not be treated as a first listing.

\par\textbf{Historical priority.} The cited paper already derives welfare conditions in its framework. No exact remaining general open-problem locator or historical priority was verified.

\subsection{Short English statement}
Estimate the specified welfare effect of a nudge using an explicit welfare criterion and identify which conclusions remain robust to uncertain preferences.

\subsection{Sources}
\begin{itemize}\raggedright
\item[\textnormal{[S1]}]\hypertarget{source:30007591:1}{} Hunt Allcott, Daniel Cohen, William Morrison, Dmitry Taubinsky. 2025. When Do "Nudges" Increase Welfare?. Abstract; AER 115(5), pp. 1555–1596. Full paper was not accessible in this audit..
\url{https://www.aeaweb.org/articles?id=10.1257/aer.20231304}.
DOI: 10.1257/aer.20231304.
{\small\textit{Supporting or current literature; see evidence qualification. The abstract explicitly presents welfare conditions involving the variance of distortions, pass-through, and optimal taxes as results. It does not establish the catalogue wording as an unresolved theorem or verify a remaining general spillover-inclusive agenda.}}
\end{itemize}
\end{document}
